\section{Appendix T -- Evidence Map and Machine-Readable Table Reference}
% R005_APPENDIX_READER_FRAME
Reader frame. This appendix supports the main manuscript by explaining source role, evidence class, audit route, and limitation boundary. It is not a detached raw dump.


\label{app:oc133-source-to-pdf-trace}

This appendix is a reader-facing evidence map. It deliberately does not reproduce the raw provenance registers, source-path digests, or machine rows inside the monograph. Those complete registers remain in the public evidence package, where they can be checked by hash. The monograph records what each evidence class does for the scientific argument and how a reader should use it.


\subsection{Public Evidence Package Contents}

The public evidence package contains the claim register, theorem registry, proof dependency graph, finite-model semantic checks, Lean source and certificate, target-blind prediction table, numeric replay QA table, domain evidence-boundary report, counterexample report, comparator and novelty registers, adversarial-review summary, and journal owner-review package index. The canonical filenames include: CLAIM\textbackslash{}\_LEDGER\textbackslash{}\_1\textbackslash{}\_3\textbackslash{}\_3.json, the corresponding public evidence-package record, FINITE\textbackslash{}\_MODEL\textbackslash{}\_CHECKS\textbackslash{}\_1\textbackslash{}\_3\textbackslash{}\_3.json, OC133V12.lean, the corresponding public evidence-package record, OC133\textbackslash{}\_TARGET\textbackslash{}\_BLIND\textbackslash{}\_PREDICTION\textbackslash{}\_TABLE.json, OC133\textbackslash{}\_NUMERIC\textbackslash{}\_REPLAY\textbackslash{}\_QA\textbackslash{}\_TABLE.json, the corresponding public evidence-package record, the corresponding public evidence-package record, OC\textbackslash{}\_1\textbackslash{}\_3\textbackslash{}\_3\textbackslash{}\_NOVELTY\textbackslash{}\_AND\textbackslash{}\_PRIORITY\textbackslash{}\_REGISTER.json, OC133\textbackslash{}\_LLM\textbackslash{}\_CERBERUS\textbackslash{}\_SUMMARY.json, SUBMISSION\textbackslash{}\_PACKAGE\textbackslash{}\_INDEX.json.


For exact SHA-256 values and source paths, use the corpus register and checksums files distributed in the release zip. Keeping those details in machine-readable artifacts avoids turning the scientific monograph into a raw hash catalogue while preserving reproducibility.


\subsection{Verification Summary}

\begin{itemize}[leftmargin=1.8em]

\item \textbf{Lean formal subset:} certificate recorded with 10 theorem references

\item \textbf{Finite semantic evidence:} 146 cases with issue count 2

\item \textbf{Bounded replay lanes:} 5 lanes with issue count 0

\item \textbf{Empirical promotion policy:} official snapshots are inputs, not empirical promotion by themselves

\item \textbf{Adversarial review:} open critical/high findings are recorded as 0/0 for the referenced review run

\item \textbf{Journal owner-review packets:} 8

\end{itemize}

\subsection{Publication-Grade Closure Rule}

A future release must fail if a primary scientific role is filled by an internal process memo, raw metadata wall, append-only surrogate, or verification package. The public manuscript must be a coherent text. Long registers may be cited and archived, but the public PDFs must teach the reader how claims, proofs, data, review, and release boundaries fit together.
